% -----------------------------------------------------------------------------
% Circuit2Code -- NeurIPS-style paper draft.
%
% The document is self-contained and compiles with a stock TeX Live install.
% To submit, replace the preamble block marked "NEURIPS SWITCH" with
%   \usepackage[final]{neurips_2024}
% and set \documentclass{article} (the official style file is not vendored here
% so that the repository builds without a network install).
% -----------------------------------------------------------------------------
\documentclass[11pt]{article}

\usepackage[letterpaper,margin=1in]{geometry}
\usepackage{times}
\usepackage[T1]{fontenc}
\usepackage[utf8]{inputenc}
\usepackage{microtype}
\usepackage{amsmath,amssymb,amsthm}
\usepackage{booktabs}
\usepackage{multirow}
\usepackage{graphicx}
\usepackage{xcolor}
\usepackage{algorithm}
\usepackage{algpseudocode}
\usepackage{hyperref}
\usepackage{url}
\usepackage{enumitem}
\usepackage{caption}
\captionsetup{font=small,labelfont=bf}

\hypersetup{colorlinks=true,linkcolor=blue!50!black,citecolor=blue!50!black,
            urlcolor=blue!50!black}

\newtheorem{theorem}{Theorem}
\newtheorem{lemma}{Lemma}
\newtheorem{definition}{Definition}
\newtheorem{proposition}{Proposition}

\newcommand{\Zp}{\mathbb{Z}_p}
\newcommand{\R}{\mathbb{R}}
\newcommand{\WQK}{W_{QK}}
\newcommand{\WOV}{W_{OV}}
\newcommand{\argmax}{\operatorname*{arg\,max}}
\newcommand{\code}[1]{\texttt{#1}}

\title{\textbf{Circuit2Code}: Verified Neuro-Symbolic Decompilation of\\
Algorithmic Transformers into Python and C}
% NOTE: anonymise for double-blind submission (replace with `Anonymous Authors`
% and drop the URL); the repository carries the full, non-anonymous history.
\author{%
  olezix\\
  \texttt{https://github.com/olezix/Circuit2Code}
}
\date{}

\begin{document}
\maketitle

% -----------------------------------------------------------------------------
\begin{abstract}
We present \textsc{Circuit2Code}, a system that decompiles a small Transformer
into an equivalent, human-readable program with a machine-checked correctness
proof: when the network implements a discrete algorithm the program is that
algorithm, and when it does not the program is the network's own trigonometric
circuit, which is exact and certified over the whole input domain. The pipeline (i) trains a
transparent attention-only Transformer on modular addition in the grokking
regime, (ii) recovers the algorithmic circuit by head ablation, direct logit
attribution and a discrete Fourier projection of the effective weight matrices
$\WQK = W_Q W_K^\top$ and $\WOV = W_V W_O^\top$, (iii) synthesizes a symbolic
program in a typed intermediate representation, (iv) proves equivalence between
the network and the program with the Z3 SMT solver over the entire finite input
domain, and (v) emits dependency-free Python and ISO C99 that compile and run
bit-exactly. Soundness is obtained by translating the network into exact rational
linear arithmetic in which the only non-algebraic quantities---the attention
probabilities---are replaced by certified outward-rounded enclosures of the
IEEE-754 forward pass, derived from Higham-style rounding-error budgets. We
distinguish three machine-checked theorems: a specification proof in quantified
Presburger arithmetic, a full-domain argmax-equivalence proof, and a bounded
logit-equivalence proof used for faithful trigonometric decompilations. On
$a+b \bmod 59$ a grokked network (100\% train, 100\% held-out) is decompiled to
the two-line modular-addition program and certified over all $3{,}481$ inputs,
while the same weights are certified as a trigonometric circuit of ten
frequencies; a partially trained and a memorising
network are certified just as rigorously but need the full basis, which turns
the qualitative ``circuits versus lookup table'' distinction into a measured
number. A $99.98\%$-correct $p=113$ network is certified over all $12{,}769$
inputs. Code, tests and this paper are reproducible from a single seed.
\end{abstract}

% -----------------------------------------------------------------------------
\section{Introduction}
\label{sec:intro}

A trained neural network is a program whose source code is a table of floating
point numbers. For models that solve discrete algorithmic tasks, this table
often implements a crisp algorithm: recent mechanistic interpretability work has
recovered Fourier multiplication circuits for modular arithmetic
\citep{nanda2023grokking}, induction heads \citep{olsson2022incontext}, and
group-operation circuits \citep{chughtai2023group}. These results are
\emph{descriptive}: they explain the network in the language of linear algebra
and features.

This paper asks a constructive question: \emph{can we turn such an explanation
into a program whose equivalence to the network is proved, not eyeballed?}
Doing so requires three things that existing interpretability tooling does not
provide: a symbolic representation with a formal semantics, a sound bridge
between floating point and real arithmetic, and a certified translation from the
symbolic representation to executable code.

\textsc{Circuit2Code} supplies all three, for the class of small algorithmic
Transformers trained on modular arithmetic. Our contributions are:

\begin{enumerate}[leftmargin=1.4em,itemsep=2pt]
  \item \textbf{A decompilation pipeline} (Section~\ref{sec:method}) that maps
        weights to a typed intermediate representation (IR), and IR to Python and
        ANSI C. The C output is C99, allocation-free, and compiles with
        \code{-Wall -Wextra -pedantic}. For a grokked $p=59$ network it is ten
        lines.
  \item \textbf{A sound float$\to$real bridge} (Section~\ref{sec:verification}).
        We compute certified enclosures of every activation with explicit
        rounding-error budgets, then convert the enclosure endpoints to exact
        rationals before they reach the solver. Attention probabilities are the
        only non-algebraic quantities; they enter the SMT query as reals
        constrained to an interval and to the probability simplex.
  \item \textbf{Three machine-checked theorems}: a quantified specification
        proof for the synthesized discrete program, a full-domain argmax
        equivalence proof over all $p^2$ inputs, and a bounded logit-equivalence
        proof for faithful trigonometric decompilations
        (Section~\ref{sec:results}).
  \item \textbf{A measured circuit-health statistic}: the fraction of spectral
        energy carried by three frequencies separates a network that has learned
        the algorithm ($97$--$99.8\%$) from one that has memorised the table
        ($27$--$43\%$), and tracks the accuracy that the decompiled circuit
        actually achieves (Section~\ref{sec:results}).
  \item \textbf{An open, reproducible implementation} with a 114-test suite in
        which the generated C is compiled and executed on every input of the
        domain, and every generated Python module is executed and compared
        against the IR's reference semantics.
\end{enumerate}

% -----------------------------------------------------------------------------
\section{Related work}
\label{sec:related}

\paragraph{Mechanistic interpretability.}
Circuit analysis decomposes a network into heads and MLP neurons and measures
their causal contribution \citep{elhage2021framework,nanda2023grokking,
olsson2022incontext,wang2022interpretability}. Attribution patching
\citep{kramar2024atp} provides a linear approximation to activation patching,
which we expose in \code{probing.py} but do not rely on for the proof.
\citet{chughtai2023group} recover group operations from learned representations;
our Fourier projection is the modular-arithmetic instance of the same idea.

\paragraph{Program synthesis from networks.}
Extracting programs from RNNs and Transformers has a long history, from
automata extraction \citep{weiss2018extracting} to differentiable program
induction. Most such methods return a program with empirical, not formal,
agreement. Causal abstraction \citep{geiger2023causal} provides the theory for when a
network implements a symbolic algorithm, and distributed alignment search
\citep{geiger2024das} turns the search for the correspondence into optimisation;
our contribution is to close the loop with an SMT proof and verified code
generation.

\paragraph{Formal verification of neural networks.}
Complete verifiers such as Marabou \citep{katz2019marabou} and $\alpha,\beta$-CROWN
\citep{zhang2018crown} prove input-output properties of piecewise-linear networks.
They target robustness specifications and do not produce programs. Conversely,
verified compilation (CompCert \citep{leroy2009compcert}) proves that a compiler
preserves semantics. We borrow the ethos of the latter---semantics preservation
with a machine-checked argument---and apply it to network$\to$program translation.

\paragraph{Grokking.}
\citet{power2022grokking} observed delayed generalisation on algorithmic
datasets; \citet{nanda2023grokking} identified the Fourier circuit responsible.
Our training setup follows theirs.

% -----------------------------------------------------------------------------
\section{Setup}
\label{sec:setup}

\subsection{Task and architecture}

Let $p$ be a prime. The dataset is the complete table
$\{(a,b)\mapsto (a+b)\bmod p : a,b\in\Zp\}$, split into a training fraction
(30\%) and a held-out fraction. An example is the three-token sequence
$[a,\,b,\,\mathrm{EQ}]$; the prediction is read from the residual stream at the
$\mathrm{EQ}$ position.

The model is a one-layer, attention-only Transformer with no LayerNorm:
\begin{align}
  x^{(0)}_t &= W_E[\text{tok}_t] \in \R^{d}, \\
  q_t, k_t, v_t &= x^{(0)}_t W_Q,\; x^{(0)}_t W_K,\; x^{(0)}_t W_V, \\
  A_{ts} &= \operatorname{softmax}_s\!\left(\tfrac{1}{\sqrt{d_h}} q_t^\top k_s\right), \\
  x^{(1)}_t &= x^{(0)}_t + \sum_{s} A_{ts}\, v_s\, W_O, \\
  \ell &= x^{(1)}_{\mathrm{EQ}}\, W_U \in \R^{p}.
\end{align}
All matrices are bias-free, so the network is a composition of linear maps and
softmaxes. This is deliberate: any bias or normalisation would have to be
translated into SMT, and LayerNorm in particular destroys the clean bilinear
structure that makes the circuit interpretable.

\subsection{Effective circuits}

Following \citet{elhage2021framework}, we summarise each head by two bilinear
forms and their token-space images:
\begin{equation}
  \WQK = W_Q W_K^\top \in \R^{d\times d}, \qquad
  \WOV = W_V W_O^\top \in \R^{d\times d},
\end{equation}
\begin{equation}
  M^{\text{att}} = W_E \WQK W_E^\top \in \R^{V\times V}, \qquad
  M^{\text{logit}} = W_E \WOV W_U \in \R^{V\times p}.
\end{equation}
$M^{\text{logit}}_{v,c}$ is the contribution of a unit attention weight on source
token $v$ to the logit of class $c$; $M^{\text{att}}_{u,v}$ is the score of key
$v$ for query $u$ (up to the constant scale).

% -----------------------------------------------------------------------------
\section{Method}
\label{sec:method}

\subsection{Circuit discovery}

For every head we (i) zero-ablate its output slice $W_O[:,h,:]$ and measure the
accuracy and logit-difference damage on the whole domain, and (ii) compute the
exact per-head, per-source-token direct logit attribution
\begin{equation}
  \ell_c \;=\; \bigl(W_E[\mathrm{EQ}]W_U\bigr)_c
  \;+\; \sum_{h}\sum_{j} A_{h,\mathrm{EQ},j}\,
        \bigl(W_E[j]\,W_{V,h}W_{O,h}W_U\bigr)_c ,
  \label{eq:dla}
\end{equation}
which holds exactly because the residual stream is linear. Heads are ranked by
zero-ablation damage; the active set is the collection of heads whose removal
costs more than a threshold.

The Fourier support of the circuit is obtained by applying the real DFT to the
token-indexed maps,
\begin{equation}
  \hat M(k) = \frac{1}{p}\sum_{x=0}^{p-1} M_{x,c}\,e^{-2\pi i kx/p},
\end{equation}
and reporting the frequencies that concentrate the energy. A grokked network is
sparse in this basis; a random network is not, which makes the statistic a
usable circuit-health signal rather than a foregone conclusion.

\subsection{Symbolic program synthesis}

The synthesizer receives the circuit specification (active heads, frequency
support, shift-invariance of $M^{\text{logit}}$) and runs two stages.

\paragraph{Stage 1: discrete motif identification.}
From a DSL of algorithmic motifs (addition, subtraction, multiplication,
projections) we instantiate every candidate $P$, evaluate it on all $p^2$
inputs, and compare with the network's argmax. A motif is accepted only if the
agreement is exact. Because the domain is finite, this check is exhaustive: it
is a proof of \emph{agreement on the domain}, not a generalisation claim, and the
SMT proof in Section~\ref{sec:verification} re-establishes it symbolically.

\paragraph{Stage 2: faithful trigonometric decompilation.}
If no motif is exact---for example for an undertrained or random network---we
emit the circuit itself. Fitting each head's score function and each column of
$M^{\text{logit}}$ with a real Fourier series over $\Zp$ gives the IR program
\begin{align}
  s_{h,j} &= \sum_{k\in K} a_{h,j,k}\cos(\omega_k j) + b_{h,j,k}\sin(\omega_k j),
  \qquad \omega_k = 2\pi k/p,\\
  P &= \operatorname{softmax}_j(s_{h,j}), \\
  \ell_c &= \beta_c + \sum_{h,j} P_{h,j}\, g_{h,c}(j), \qquad
  \text{answer} = \argmax_c \ell_c ,
\end{align}
with $g_{h,c}$ the truncated Fourier series of $M^{\text{logit}}_{\cdot,c}$. The
truncation error is measured and reported; with $|K|$ frequencies the
reconstruction is exact whenever $2|K|+1 \ge p$.

\subsection{IR and code generation}

The IR (\code{ast\_nodes.py}) is a small expression language with three
evaluators: NumPy reference semantics, outward-rounded interval semantics, and
two source emitters. Nodes cover the algorithmic primitives of interest
(\code{Mod}, \code{Add}, \code{Sin}, \code{Cos}, \code{Linear}, \code{Threshold},
\code{ArgMax}, \code{IfElse}) plus \code{Softmax} and its projections. SymPy is
used only as a *conservative* simplifier: a rewrite is kept only if it is
numerically equivalent on a fixed grid.

The Python emitter produces a self-contained module with a docstring, a stable
softmax/argmax implementation, and a ``\code{if \_\_name\_\_ == '\_\_main\_\_'}'' demo.
The C emitter produces ISO C99: integer types when the program is integer-valued
(bit-exact), \code{double} otherwise, static tables for linear maps, and no heap
allocation, so there is nothing to leak.

% -----------------------------------------------------------------------------
\section{Formal verification}
\label{sec:verification}

\subsection{From floats to reals}

\begin{definition}[Certified enclosure]
For an activation $x$ computed by the IEEE-754 forward pass, an interval
$[\underline{x},\overline{x}]$ is a certified enclosure if the computed value
provably lies inside it, using only the semantics of IEEE-754 round-to-nearest
and outward rounding, with no reference to the value being checked.
\end{definition}

We propagate enclosures through the network. For a dot product of length $n$,
Higham's standard bound gives
\begin{equation}
  \bigl|\operatorname{fl}(\langle x,w\rangle) - \langle x,w\rangle\bigr|
  \le \gamma_n \sum_i |x_i w_i|, \qquad
  \gamma_n = \frac{n u}{1-nu}, \quad u = 2^{-53}.
\end{equation}
Applying this to every matmul, an interval softmax (monotonicity of $\exp$ plus
outward-rounded division and the simplex identity $\sum_j P_j = 1$), and
outward-rounded residual additions yields enclosures whose width is at most
$\sim 10^{-13}$ for the models in this paper---twelve orders of magnitude
tighter than the logit margins of a trained network, but computed without ever
looking at the answer.

\subsection{SMT encodings}

\paragraph{Specification (exact, integer).}
Wrap the IR program in a quantified Presburger formula:
\begin{equation}
  \forall a,b \in [0,p): \bigl(P(a,b) = a+b \wedge P(a,b) < p\bigr)
  \;\vee\; \bigl(P(a,b) = a+b-p \wedge P(a,b) \ge 0\bigr).
\end{equation}
No floating point appears. The negation is handed to Z3; \code{unsat} means the
synthesized program \emph{is} modular addition, independent of the network.

\paragraph{Model equivalence (rational, full domain).}
For each input $x=(a,b)$ we introduce one real variable $P_{h,j}$ per attention
probability, constrain $0 \le P_{h,j} \le 1$, $\sum_j P_{h,j} = 1$, and
$P_{h,j} \in [\underline{P}_{h,j}, \overline{P}_{h,j}]$ (the certified
enclosure), and write the logits as exact rational linear forms in these
variables using the weights interpreted as exact binary rationals. The query is
\begin{equation}
  \bigvee_{x \in \Zp^2} \Bigl(\text{enclosures}(x) \;\wedge\;
  \bigvee_{c \ne P(x)} \ell_c(x) \ge \ell_{P(x)}(x)\Bigr).
  \label{eq:query}
\end{equation}
Classes with $\overline{\ell}_c < \underline{\ell}_{P(x)}$ cannot satisfy the
inner disjunction (a sound, exact pruning of *classes*, never of inputs). The
negation of \eqref{eq:query} is what we prove.

\begin{theorem}[Soundness of the equivalence query]
\label{thm:sound}
If the certified enclosures contain the forward-pass values (Definition 1) and
\eqref{eq:query} is unsatisfiable, then for every input in $\Zp^2$ the network's
argmax equals the program's output.
\end{theorem}
\begin{proof}[Proof sketch]
Suppose some input $x^\star$ had a different argmax, realised by the true
attention assignment $P^\star = P(x^\star)$ and a class $c^\star \ne P(x^\star)$
with $\ell_{c^\star}(x^\star) \ge \ell_{P(x^\star)}(x^\star)$. By Definition 1,
$P^\star$ satisfies the enclosure and simplex constraints; by exact rational
arithmetic, the linear forms evaluate to the true logits. Then
$(x^\star, P^\star)$ is a satisfying assignment of \eqref{eq:query}, contradicting
unsatisfiability.
\end{proof}

\paragraph{Bounded logit equivalence.}
For faithful trig decompilations the target is numeric:
$|\ell_c^{\text{net}}(x) - \ell_c^{\text{prog}}(x)| \le \varepsilon$ for all $x$
and $c$, encoded by declaring each network logit inside its enclosure and asking
for a deviation beyond $\varepsilon$.

\paragraph{Robustness.}
Widening every enclosure by $\pm\delta$ yields a strictly stronger theorem: the
argmax is invariant to a $\delta$-perturbation of every attention probability.
This is also the regime in which the linear SMT solver performs real work, since
the enclosure-based pruning no longer decides the query.

\subsection{Counterexample-guided refinement}

If the query is satisfiable, the witness is replayed through the real network.
A genuine mismatch triggers refinement: the synthesizer is asked for the
faithful trigonometric circuit, which is then verified with the bounded-logit
theorem. A witness that only touches the edge of an enclosure is reported as
\code{unknown}---never as a proof and never as a counterexample.

% -----------------------------------------------------------------------------
% BEGIN GENERATED: experiments  (scripts/render_report.py)
\section{Experiments}
\label{sec:results}

\subsection{Setup}

All models are one-layer, attention-only Transformers trained full-batch with AdamW ($\beta=(0.9,0.98)$, weight decay $1.0$) on one CPU core, single BLAS thread (Table~\ref{tab:checkpoints}); the $p=13$ and $p=23$ instances use $d=64$ with four heads, all others $d=128$. The learning rate is half-cosine annealed from $10^{-2}$ to $2\%$ of its value, except for the memoriser, which uses a constant $10^{-3}$. Runs are deterministic given the seed \emph{and} the library versions, which \code{artifacts/environment.json} records; everything reported here was measured with Python 3.11.2, PyTorch 2.15.0+cu132, NumPy 2.4.6 and Z3 5.3.3. The checkpoints themselves are not distributed: \code{make train} rebuilds them and \code{make results} regenerates every number in this section from the JSON artefacts.

\begin{table}[t]
\centering
\small
\begin{tabular}{@{}lrrrrrr@{}}
\toprule
Checkpoint & $p$ & train frac. & params & epochs & held-out & full domain \\
\midrule
G (grokked) & 59 & 0.70 & 80{,}768 & 3{,}250 & 1.000 & 1.000 \\
T (partial) & 59 & 0.30 & 80{,}768 & 30{,}000 & 0.7739 & 0.8417 \\
M (memoriser) & 59 & 0.30 & 80{,}768 & 40{,}000 & 0.3028 & 0.5119 \\
$p=13$ & 13 & 0.50 & 18{,}112 & 60{,}000 & 0.7294 & 0.8639 \\
$p=23$ & 23 & 0.50 & 19{,}392 & 30{,}000 & 0.7434 & 0.8715 \\
$p=113$ & 113 & 0.70 & 94{,}592 & 20{,}000 & 0.9992 & 0.9998 \\
\bottomrule
\end{tabular}
\caption{The checkpoints of this section. ``held-out'' is the fraction of the domain the optimiser never saw, evaluated with the seed the model was trained with; ``full domain'' scores all $p^2$ inputs.}\label{tab:checkpoints}
\end{table}

\subsection{Delayed generalisation and circuit structure}

Table~\ref{tab:training} is the trajectory of checkpoint T, trained on $30\%$ of the domain. Training accuracy reaches $100\%$ early while held-out accuracy sits at chance ($1/59\approx1.7\%$); held-out accuracy only starts to move at epoch 14,750, when the annealed learning rate is small, and the run ends at $0.774$ --- delayed generalisation that never reaches exactness.

\begin{table}[t]
\centering
\small
\begin{tabular}{@{}rrrr@{}}
\toprule
Epoch & Train acc. & Test acc. & Test loss \\
\midrule
2{,}500 & 1.000 & 0.303 & 11.06 \\
10{,}000 & 1.000 & 0.364 & 6.23 \\
12{,}000 & 1.000 & 0.367 & 7.84 \\
15{,}000 & 1.000 & 0.572 & 3.52 \\
18{,}000 & 1.000 & 0.638 & 2.39 \\
22{,}500 & 1.000 & 0.616 & 4.05 \\
30{,}000 & 1.000 & 0.774 & 2.07 \\
\bottomrule
\end{tabular}
\caption{The delayed-generalisation trajectory of checkpoint T ($30\%$ of the domain): held-out accuracy is still at chance after thousands of epochs of perfect training accuracy.}\label{tab:training}
\end{table}

Checkpoint G, trained on $70\%$ of the domain with the same schedule, is at chance until epoch 1,250 and then transitions to exact generalisation in a few hundred epochs; training stops 2,000 epochs after the transition, at epoch 3,250, when the network is exact on the training set, on the held-out set and on the whole domain. The transition epoch is a property of the numerical environment --- we have seen the same code and seed grok anywhere between $1{,}250$ and $16{,}750$ epochs in different PyTorch builds --- whereas the circuit recovered from the endpoint has been identical every time.

\begin{table}[t]
\centering
\small
\begin{tabular}{@{}lrrrr@{}}
\toprule
Head & $\Delta$acc.\ (zero-ablation) & $\Delta$logit & top frequencies & top-3 energy \\
\midrule
\multicolumn{5}{@{}l}{\emph{checkpoint G (grokked)}, held-out $1.000$}\\
L0H0 & $0.872$ & $39.80$ & $22,\,15,\,28$ & $0.995$ \\
L0H1 & $0.888$ & $33.50$ & $22,\,15,\,28$ & $0.994$ \\
L0H2 & $0.834$ & $17.03$ & $22,\,28,\,15$ & $0.982$ \\
L0H3 & $0.757$ & $17.00$ & $28,\,22,\,3$ & $0.979$ \\
\midrule
\multicolumn{5}{@{}l}{\emph{checkpoint T (partial)}, held-out $0.774$}\\
L0H0 & $0.598$ & $4.21$ & $6,\,29,\,11$ & $0.928$ \\
L0H1 & $0.649$ & $28.86$ & $29,\,11,\,6$ & $0.958$ \\
L0H2 & $0.696$ & $6.00$ & $11,\,29,\,22$ & $0.955$ \\
L0H3 & $0.698$ & $25.26$ & $11,\,29,\,6$ & $0.951$ \\
\midrule
\multicolumn{5}{@{}l}{\emph{checkpoint M (memoriser)}, held-out $0.303$}\\
L0H0 & $0.380$ & $3.16$ & $16,\,25,\,6$ & $0.324$ \\
L0H1 & $0.379$ & $3.17$ & $16,\,20,\,2$ & $0.287$ \\
L0H2 & $0.374$ & $2.66$ & $20,\,11,\,6$ & $0.241$ \\
L0H3 & $0.375$ & $2.56$ & $16,\,2,\,4$ & $0.373$ \\
\bottomrule
\end{tabular}
\caption{Recovered circuits. ``Top-3 energy'' is the fraction of the spectral energy of $M^{\text{logit}}$ carried by the three dominant frequencies; ``$\Delta$acc.'' is the accuracy lost when the head's output projection is zeroed on the full domain.}\label{tab:circuit}
\end{table}

Three statistics separate the checkpoints (Table~\ref{tab:circuit}).

\emph{Spectral concentration.} In checkpoint G three frequencies carry $97.9$--$99.5\%$ of the energy of every head's token$\to$logit map, and the same modes dominate all four heads. In checkpoint T the concentration falls to $92.8$--$95.8\%$ and the dominant set differs from head to head; in the memoriser the three leading frequencies carry only $24.1$--$37.3\%$, because the map is a lookup table spread over the whole basis.

\emph{Causal load.} Zero-ablating any single head's output projection costs $75.7$--$88.8\%$ accuracy in checkpoint G and $59.8$--$69.8\%$ in checkpoint T, so every head takes part in the computation, but only $37.4$--$38.0\%$ in the memoriser, where the work is replicated across heads.

\emph{Circuit size.} Table~\ref{tab:size} rebuilds the circuit from its $|K|$ strongest frequencies. For checkpoint G, \emph{10} of the 29 available frequencies already reproduce the network's argmax on every input, and Z3 proves the full-domain equivalence of that program; the exact reconstruction of the \emph{logits} needs the whole basis, whose truncation error is $9.17\times10^{-13}$.

\begin{table}[t]
\centering
\small
\begin{tabular}{@{}rrrr@{}}
\toprule
$|K|$ & agreement & logit error & \code{model\_equivalence} \\
\midrule
1 & $0.2390$ & $6.05\times10^{1}$ & refuted (by counting) \\
2 & $0.5248$ & $3.14\times10^{1}$ & refuted (by counting) \\
3 & $0.8995$ & $1.40\times10^{1}$ & refuted (by counting) \\
4 & $0.9879$ & $1.05\times10^{1}$ & refuted (by counting) \\
5 & $0.9948$ & $1.03\times10^{1}$ & refuted (by counting) \\
6 & $0.9954$ & $8.08\times10^{0}$ & refuted (by counting) \\
7 & $0.9989$ & $8.10\times10^{0}$ & refuted (by counting) \\
8 & $0.9997$ & $8.31\times10^{0}$ & refuted (by counting) \\
9 & $0.9994$ & $7.66\times10^{0}$ & refuted (by counting) \\
10 & $\mathbf{1.0000}$ & $6.66\times10^{0}$ & proved \\
11 & $0.9994$ & $5.68\times10^{0}$ & refuted (by counting) \\
12 & $0.9997$ & $4.88\times10^{0}$ & refuted (by counting) \\
29 & $\mathbf{1.0000}$ & $9.17\times10^{-13}$ & proved \\
\bottomrule
\end{tabular}
\caption{Full-domain argmax agreement between a $|K|$-frequency truncation of the faithful circuit and checkpoint G. The smallest exact truncation is $|K|=10$ and its equivalence is \emph{proved}; every row with agreement below one is refuted by counting, so no SMT query is needed. The curve is not monotone --- a truncation is a re-projection, not a nested approximation of the argmax.}\label{tab:size}
\end{table}

\subsection{Decompilation and proof}

\paragraph{A discrete program.} For checkpoint G the synthesizer's first stage succeeds: the argmax of the network is $a+b \bmod p$ on all 3,481 inputs, so the emitted program is

\begin{center}
\code{def add\_mod\_p(a, b): return (a + b) \% 59}
\end{center}

(4 IR nodes). Z3 discharges the specification --- the quantified Presburger statement $\forall a,b\in\Zp: (P=a+b \wedge P<p) \vee (P=a+b-p \wedge P\ge 0)$ --- in 0.02\,s and the full-domain equivalence with the network in 1.42\,s (Table~\ref{tab:verification}). The emitted C is a 35-line file, allocation-free, and compiles with \code{-Werror} and no diagnostics (Table~\ref{tab:export}).

\paragraph{The same network as a trigonometric circuit.} Building the faithful circuit from checkpoint G's weights reproduces its argmax on all 3,481 inputs at the full basis and, as Table~\ref{tab:size} shows, already at $|K|=10$; Z3 proves that program equivalent to the network over the whole domain. The logit-level statement is \emph{not} true at that truncation --- the logits differ by up to 6.66 in absolute value --- which is exactly why the prover is asked for the argmax statement: the same network is simultaneously the discrete algorithm and an approximate trigonometric circuit, and both readings are machine-checked.

\paragraph{Networks without a motif.} For checkpoints T and M stage 1 finds no exact motif (the faithful circuit agrees with the network on $0.842$ and $0.512$ of the domain, and the residual is a genuine error of the trained weights, not of the extraction), so the faithful circuit is emitted and proved equivalent over the whole domain. Certification does not depend on the network being right, only on it being a well-defined function. A randomly initialised $p=11$ network --- the worst case for a decompiler --- is handled the same way: full basis, truncation error $5.33\times10^{-15}$ on the scores and $2.49\times10^{-14}$ on the logits, equivalence proved.

\paragraph{A second modulus.} The same pipeline, unchanged, on $p=113$: a 94,592-parameter network trained on $70\%$ of the 12,769 pairs reaches 0.99977 of the domain (3 inputs wrong). Stage 1 declines to claim a motif, the faithful circuit (11,055 nodes) is emitted, and Z3 certifies both the argmax statement over all 12,769 inputs in 501\,s and the logit statement, with a deployment gap of $1.83\times10^{-4}$.

\begin{table}[t]
\centering
\small
\begin{tabular}{@{}llll@{}}
\toprule
Target & Theorem & Domain & Result \\
\midrule
$p=59$ G & \code{spec} & 3{,}481 inputs & proved, 0.02\,s \\
$p=59$ G & \code{model\_equivalence} & 3{,}481 inputs & proved \\
$p=59$ T & \code{model\_equivalence} & 3{,}481 inputs & proved \\
$p=59$ T & \code{logit\_equivalence} & 3{,}481 inputs & proved \\
$p=59$ M & \code{model\_equivalence} & 3{,}481 inputs & proved \\
$p=59$ M & \code{logit\_equivalence} & 3{,}481 inputs & proved \\
$p=113$ & \code{model\_equivalence} & 12{,}769 inputs & proved \\
$p=113$ & \code{logit\_equivalence} & 12{,}769 inputs & proved \\
$p=13$ & \code{model\_equivalence} & 169 inputs & proved \\
$p=13$ & \code{logit\_equivalence} & 169 inputs & proved \\
$p=23$ & \code{model\_equivalence} & 529 inputs & proved \\
$p=23$ & \code{logit\_equivalence} & 529 inputs & proved \\
$p=11$ random & \code{model\_equivalence} & 121 inputs & proved \\
$p=11$ random & \code{logit\_equivalence} & 121 inputs & proved \\
\bottomrule
\end{tabular}
\caption{Proof obligations discharged by Z3. Every query covers the full domain; no input is sampled or skipped. \code{spec} is the quantified integer specification, \code{model\_equivalence} the full-domain argmax statement and \code{logit\_equivalence} the bounded-logit statement used when no discrete program exists.}\label{tab:verification}
\end{table}

\subsection{Robustness: where the solver does real work}

For most proved queries the certified enclosures are tight enough that every input is decided by exact rational comparison and Z3 only confirms unsatisfiability. Widening the enclosures by $\pm\delta$ turns the theorem into a strictly stronger robustness claim and forces the solver to reason inside the simplex (Table~\ref{tab:robustness}).

\begin{table}[t]
\centering
\small
\begin{tabular}{@{}rlr@{}}
\toprule
$\delta$ & verdict & inputs decided by Z3 \\
\midrule
$0$ & proved & $0/169$ \\
$0.001$ & proved & $6/169$ \\
$0.01$ & refuted & n/a \\
\bottomrule
\end{tabular}
\caption{Robustness of the $p=13$ circuit to a uniform perturbation of every attention probability. ``inputs decided by Z3'' counts the inputs whose winning margin could no longer be certified by exact rational comparison. A witness inside the perturbation set is reported as \code{refuted}, not \code{unknown}: for a robustness statement it is a counterexample, even though the unperturbed network still agrees with its program on that input.}\label{tab:robustness}
\end{table}

\subsection{Cost}

\begin{itemize}[leftmargin=1.4em,itemsep=2pt]
  \item \emph{Discovery and synthesis.} Circuit discovery takes under a second; synthesis takes 0.21\,s for the grokked checkpoint and 6.7\,s when the faithful $p=59$ circuit (5,763 nodes) is emitted.
  \item \emph{Proof.} 1.42\,s for the grokked checkpoint and 41.2\,s for the faithful $p=59$ circuit, dominated by building the certified enclosures rather than by Z3; 501\,s for $p=113$, whose domain is 3.7$\times$ larger.
  \item \emph{Code generation.} Under a second; compiling the 2.08\,MB source of checkpoint T takes 2.50\,s.
  \item \emph{Enclosure width.} With two-ulp outward padding per operation the logit enclosures stay within $1.5\times10^{-14}$ of the float64 forward pass, which the test suite validates for every input.
  \item \emph{Deployment gap.} float32 and float64 logits differ by at most 6.25e-5 ($p=59$), 1.11e-5 ($p=13$) and 4.25e-5 ($p=23$) in absolute value. The proof is about the float64 network and this gap is reported, never assumed away.
  \item \emph{Training.} Dominates everything: 5 minutes on one CPU core for checkpoint G, 18 minutes for checkpoint T.
\end{itemize}

\begin{table}[t]
\centering
\small
\begin{tabular}{@{}lr@{}}
\toprule
Measurement & Value \\
\midrule
IR nodes & 5{,}763 \\
Python source & 2.21\,MB \\
C source & 2.08\,MB, 869 lines \\
\code{gcc} time & 2.50\,s \\
\code{gcc} diagnostics & 0 \\
input domain & 3{,}481 pairs \\
C vs.\ Python, all inputs & max $|\Delta| = 0.0$ (bit-identical) \\
heap allocation & none \\
\bottomrule
\end{tabular}
\caption{The emitted code of checkpoint T, the largest artefact, compiled with \code{-std=c99 -O3 -Wall -Wextra -pedantic -Werror}. The C and the Python agree on the whole domain, and at 599 bytes per input the C cannot be a lookup table in disguise.}\label{tab:export}
\end{table}

\subsection{What is not claimed}
\label{sec:notclaimed}

Checkpoint G was trained on $70\%$ of the domain, and it is exact on the held-out $30\%$ as well. The scarcer-data setting does not work here: checkpoint T, trained on $30\%$ with the same schedule, plateaus at $0.774$ held-out accuracy, and we report the fraction that works rather than presenting a hand-picked success as the typical behaviour. Nothing in the method depends on which of the two succeeds --- the decompiler is exact and proved in both cases --- and the difference shows up quantitatively as circuit size, $|K|=10$ versus the whole basis.

The verified object is the float64 forward pass; the gap to the float32 model that would be deployed is measured over the whole domain and reported, never assumed away. Softmax is not an SMT primitive, so its semantics are carried by the enclosures rather than by the solver. The proof quantifies over the inputs and over the certified enclosures, not over all possible attention distributions. Enclosures are sound by construction, but they are machinery, not a proof assistant: the trust base is our interval arithmetic plus Z3. Only finite, enumerable domains can be certified, and only one-layer attention-only models can currently be decompiled.
% END GENERATED: experiments

% -----------------------------------------------------------------------------
\section{Discussion and limitations}
\label{sec:discussion}

\paragraph{What the proof does and does not say.}
The theorem quantifies over the finite domain of inputs and over the certified
enclosures of the attention probabilities; it does not quantify over attention
distributions the network cannot produce. The verified object is the float64
forward pass; the float32 deployment gap is measured (order $10^{-5}$ absolute on
logits) and reported alongside every proof. Softmax is not an SMT primitive, so
its semantics are carried by the enclosures rather than by the solver; a fully
symbolic treatment of softmax would require transcendental-function reasoning
that current SMT solvers do not provide.

\paragraph{Scope.}
The symbolic encoder and faithful circuit target one-layer attention-only
models. The enclosure engine and the SMT encoder already handle several layers
and a ReLU MLP, so a deeper model can be *proved* against a given program, but
no decompiler is provided for it, and we claim no such experiment. Motif
identification is exhaustive over a small DSL; programs outside the DSL fall
back to the faithful circuit rather than to open-ended synthesis.

\paragraph{Tasks other than modular arithmetic.}
We do not report Dyck-1 bracket balancing, which we had planned as a second
domain. The reason is structural rather than incidental: our architecture has no
positional embeddings, so it is permutation-invariant and cannot represent a
prefix-dependent language at all. Dyck support needs learned positional
embeddings, which change the forward pass and therefore the SMT encoder. The
*symbolic* half is already expressible in the IR --- a word is balanced iff
every prefix sum is non-negative and the total is zero --- but without a model
that can learn the task there is nothing to decompile.

\paragraph{Broader impact.}
Turning a network into auditable code is a step towards accountability for
deployed models, and a proof-carrying artefact is easier to review than a weight
matrix. The same machinery could be misused to certify a circuit that implements
a harmful algorithm; the method is agnostic to the task semantics, which is
precisely why the specification step is kept explicit and human-readable.

% -----------------------------------------------------------------------------
\section{Conclusion}
\label{sec:conclusion}

We have shown that a Transformer trained on modular arithmetic can be
decompiled into readable Python and allocation-free C, with the equivalence
between network and program proved by an SMT solver over the whole input domain:
into the two-line algorithm when the network implements it, and into its
trigonometric circuit otherwise. The key technical ingredient is a sound and
tight bridge between IEEE-754 arithmetic and the exact rationals that the solver
understands. The measured spectral statistic --- $97.9$--$99.5\%$ of the energy of
a grokked network in three frequencies, versus $24.1$--$37.3\%$ for a memoriser ---
together with the ten-versus-$29$ frequency gap in Table~\ref{tab:size}, is a
circuit-health signal that future work can use to decide when an extracted
program is the \emph{algorithm} rather than a description of one network. The
pipeline itself carries over to a second modulus without changes: a
$99.98\%$-correct $p=113$ network is decompiled and certified over all $12{,}769$
of its inputs. We release
the implementation, the tests and the generation pipeline so that every claim
can be re-derived end to end from a single seed.

% -----------------------------------------------------------------------------
\section*{Reproducibility statement}
All experiments are produced by the commands in the accompanying repository.
Every model and every data split is built from an explicit seed (42 for the unit
tests, 3 for the reported runs), training is full-batch with a pinned BLAS thread
count, and a rerun in the same software environment is bitwise identical (a unit
test asserts this). Across PyTorch releases the \emph{trajectory} can differ
substantially (the same seed has grokked anywhere between 1{,}250 and 16{,}750
epochs in different PyTorch builds), which is why every claim we make about a
network is a certified
property of its weights, not a statement about an epoch number. Each command
writes a JSON artefact: the checkpoint, circuit and proof rows of
Tables~\ref{tab:checkpoints}--\ref{tab:verification} come from
\code{artifacts/paper\_tables.json}, the truncation sweep of
Table~\ref{tab:size} from \code{artifacts/truncation\_sweep\_p59.json}, the
compile measurements of Table~\ref{tab:export} from
\code{artifacts/export\_measurements.json}, and the trajectory of
Table~\ref{tab:training} is the \code{<checkpoint>.history.json} that
\code{train} writes next to each checkpoint.
The $p=113$ checkpoint of Section~\ref{sec:results} comes from

\begin{quote}\small
\code{python demo/cli.py train --p 113 --d-model 128 --n-heads 4 --epochs 20000
--lr 1e-2 --lr-schedule cosine --weight-decay 1.0 --train-fraction 0.7 --seed 3
--torch-threads 2 --checkpoint-every 2000 --out artifacts/checkpoints/g113.pt}
\end{quote}

\noindent (two BLAS threads on one CPU, 1h37m for 20{,}000 epochs), after which
\code{verify}
spends 501\,s on the 12{,}769-input proof.

% -----------------------------------------------------------------------------
\begin{thebibliography}{99}
\small

\bibitem{chughtai2023group}
B.~Chughtai, L.~Chan, and N.~Nanda.
\newblock A toy model of universality: Reverse engineering how networks learn
  group operations.
\newblock \emph{arXiv:2302.03025}, 2023.

\bibitem{elhage2021framework}
N.~Elhage, N.~Nanda, C.~Olsson, et al.
\newblock A mathematical framework for Transformer circuits.
\newblock \emph{Transformer Circuits Thread}, 2021.

\bibitem{katz2019marabou}
G.~Katz, D.~A. Huang, D.~Ibeling, et al.
\newblock The Marabou framework for verification and analysis of deep neural
  networks.
\newblock In \emph{Computer Aided Verification (CAV)}, 2019.

\bibitem{kramar2024atp}
J.~Kram\'ar, T.~Lieberum, R.~Shah, and N.~Nanda.
\newblock AtP*: An efficient and scalable method for localizing LLM behaviour to
  components.
\newblock \emph{arXiv:2403.00745}, 2024.

\bibitem{leroy2009compcert}
X.~Leroy.
\newblock Formal verification of a realistic compiler.
\newblock \emph{Communications of the ACM}, 52(7):107--115, 2009.

\bibitem{nanda2023grokking}
N.~Nanda, L.~Chan, T.~Lieberum, J.~Smith, and J.~Steinhardt.
\newblock Progress measures for grokking via mechanistic interpretability.
\newblock In \emph{International Conference on Learning Representations (ICLR)},
  2023.

\bibitem{olsson2022incontext}
C.~Olsson, N.~Elhage, N.~Nanda, et al.
\newblock In-context learning and induction heads.
\newblock \emph{Transformer Circuits Thread}, 2022.

\bibitem{power2022grokking}
A.~Power, Y.~Burda, H.~Edwards, I.~Babuschkin, and V.~Misra.
\newblock Grokking: Generalization beyond overfitting on small algorithmic
  datasets.
\newblock \emph{arXiv:2201.02177}, 2022.

\bibitem{wang2022interpretability}
K.~Wang, A.~Variengien, A.~Conmy, B.~Shlegeris, and J.~Steinhardt.
\newblock Interpretability in the wild: a circuit for indirect object
  identification in GPT-2 small.
\newblock In \emph{ICLR}, 2023.

\bibitem{weiss2018extracting}
G.~Weiss, Y.~Goldberg, and E.~Yahav.
\newblock Extracting automata from recurrent neural networks using queries and
  counterexamples.
\newblock In \emph{International Conference on Machine Learning (ICML)}, 2018.

\bibitem{geiger2023causal}
A.~Geiger, D.~Ibeling, A.~Zur, M.~Chaudhary, S.~Chauhan, J.~Huang, A.~Arora,
  Z.~Wu, N.~Goodman, C.~Potts, and T.~Icard.
\newblock Causal abstraction: A theoretical foundation for mechanistic
  interpretability.
\newblock \emph{arXiv:2301.04709}, 2023 (v4, 2025).

\bibitem{geiger2024das}
A.~Geiger, Z.~Wu, C.~Potts, T.~Icard, and N.~Goodman.
\newblock Finding alignments between interpretable causal variables and
  distributed neural representations.
\newblock In \emph{Proceedings of the 3rd Conference on Causal Learning and
  Reasoning (CLeaR)}, 2024.

\bibitem{zhang2018crown}
H.~Zhang, T.-W. Weng, P.-Y. Chen, C.-J. Hsieh, and L.~Daniel.
\newblock Efficient neural network robustness certification with general
  activation functions.
\newblock In \emph{NeurIPS}, 2018.

\end{thebibliography}

% -----------------------------------------------------------------------------
\appendix
\section{Implementation details}
\label{app:impl}

\paragraph{Enclosure implementation.}
\code{compute\_logit\_bounds} propagates an (\code{lo}, \code{hi}) pair per
activation. Matmuls use the corner method plus a $\gamma_n$ error budget;
softmax shifts by the upper bound of the maximum score (making the exponential
well conditioned) and computes monotone exponential bounds padded by two ulps;
each probability is bounded by the ratio of an extreme numerator and an extreme
denominator, which is sound by construction. All padding uses
\code{numpy.nextafter}.

\paragraph{Z3 encoding of the IR.}
Integer programs are translated with Z3's Euclidean \code{Mod}; real programs use
exact rational constants built from \code{fractions.Fraction}. Nodes without an
SMT counterpart (\code{Sin}, \code{Cos}, \code{Softmax}, \code{ArgMax}) raise
\code{NotImplementedError}, and the verifier falls back to the enclosure path
instead of silently approximating.

\paragraph{Batching.}
The equivalence query is split into batches (\code{VerifyConfig.batch\_size},
default 128). Each batch is a single \code{SolverFor("QF\_LRA")} call, so the
solver sees the disjunction over thousands of inputs at once.

\paragraph{Testing.}
The suite contains 118 tests, including an end-to-end test of the CLI, one test
per measurement script, and one test that re-derives the paper's headline claim
from the grokked checkpoint (skipped when the checkpoint is absent). Codegen
tests execute the emitted Python and compile the emitted C with \code{-Werror}, then compare against the IR on the
complete domain. Verifier tests check enclosure soundness against the float64
forward pass on every input, and check that a deliberately wrong program is
refuted with a genuine counterexample.

\section{IR grammar}
\label{app:grammar}

\begin{align*}
  e &::= c \mid x \mid -e \mid e \pm e \mid e \cdot e \mid e / e
       \mid e \bmod m \mid \sin e \mid \cos e \\
    &\quad\mid \textstyle\sum_i e_i \mid W e \mid \mathrm{softmax}(e_1,\dots,e_n)
       \mid \mathrm{softmax}(\cdot)[i] \\
    &\quad\mid \mathrm{step}(e - t) \mid \argmax_i(e_1,\dots,e_n)
       \mid \mathrm{if}\ b\ \mathrm{then}\ e_1\ \mathrm{else}\ e_2 ,
\end{align*}
with $b$ a comparison or boolean combination thereof. Programs are pairs
$(x_1,\dots,x_n \mapsto e)$ plus the modulus of the ambient ring.

\end{document}
